\begin{figure}[H]
    \centering
    \caption{Paradoxo do Barbeiro}
    \label{fig:paradoxo_barbeiro_sagital}

    \begin{tikzpicture}[
            scale=0.8,
            transform shape,
            font=\sffamily,
            >=stealth
        ]

        \def\ColorSet{\ifbool{darkmode}{gray!50}{gray!60}}
        \def\ColorA{\ifbool{darkmode}{DarkModeLink}{LightModeLink}}
        \def\ColorB{\ifbool{darkmode}{DarkModeRed}{LightModeRed}}
        \def\ColorParadox{\ifbool{darkmode}{DarkModeViolet}{LightModeViolet}}
        \def\ColorText{\ifbool{darkmode}{DarkModeText}{black}}

        \draw[thick, \ColorSet, fill=\ColorSet, fill opacity=0.05] (0,0) ellipse (5cm and 3.5cm);
        \node[\ColorText, font=\Large\bfseries] at (0, 4.2) {Homens da Cidade};

        \draw[thick, \ColorSet, dashed] (0, -3.5) -- (0, 3.5);

        \node[\ColorA, font=\small\bfseries, align=center] at (-2, 2.5) {Se barbeiam};
        \node[\ColorB, font=\small\bfseries, align=center] at (1.9, 2.5) {Não se barbeiam};

        \node[circle, fill=\ColorA, inner sep=1.5pt, label={[text=\ColorA]left:\(h_1\)}] (h1) at (-2.5, 1) {};
        \node[circle, fill=\ColorA, inner sep=1.5pt, label={[text=\ColorA]left:\(h_2\)}] (h2) at (-2.5, -1.5) {};

        \node[circle, fill=\ColorB, inner sep=1.5pt, label={[text=\ColorB]right:\(h_3\)}] (h3) at (2.5, 1) {};
        \node[circle, fill=\ColorB, inner sep=1.5pt, label={[text=\ColorB]right:\(h_4\)}] (h4) at (2.5, -1.5) {};

        \node[circle, fill=\ColorParadox, inner sep=2pt, label={[text=\ColorParadox, font=\bfseries]left:\(b\)}] (b) at (0, -0.5) {};

        \draw[->, \ColorA, thick] (h1) to[out=135, in=45, distance=1.2cm] node[above, font=\footnotesize] {barbeia} (h1);
        \draw[->, \ColorA, thick] (h2) to[out=135, in=45, distance=1.2cm] node[above, font=\footnotesize] {barbeia} (h2);

        \draw[->, \ColorParadox, thick] (b) -- (h3) node[midway, sloped, above, font=\footnotesize] {barbeia};
        \draw[->, \ColorParadox, thick] (b) -- (h4) node[midway, sloped, below, font=\footnotesize] {barbeia};

        \draw[->, \ColorParadox, ultra thick, dashed] (b) to[out=60, in=120, distance=2cm] node[above, font=\Large\bfseries] {?} (b);

        % \node[\ColorText, align=right, font=\footnotesize, text width=4.5cm] at (-2.5, -4.8) {
        %     Se a seta existir, \(b\) deveria estar no lado esquerdo (mas ele não pode barbear quem se barbeia).
        % };

        % \node[\ColorText, align=left, font=\footnotesize, text width=4.5cm] at (2.5, -4.8) {
        %     Se a seta não existir, \(b\) pertence ao lado direito (mas a regra o obriga a barbear todos deste lado).
        % };

    \end{tikzpicture}

    \fonte{Elaborado pelo autor (2026).}
\end{figure}